\section{Difesa
dell\textquotesingle Architettura}\label{difesa-dellarchitettura}

Spunti:

\begin{itemize}
\tightlist
\item
  perche\textquotesingle{} Debian 13?
\item
  perche\textquotesingle{} ZFS?
\item
  perche\textquotesingle{} OVN in verifica?
\item
  perche\textquotesingle{} Birdge L2 in produzione?
\item
  perche\textquotesingle{} Incus?
\end{itemize}

Codice:

\begin{itemize}
\tightlist
\item
  Motivazione per Async Context Manager
\item
  perche ContainerEngine ed IncusEngine?
\item
  perche esiste il problema della concorrenza?
\item
  perche disaccoppiare la creazione dei comandi?
\item
  perche gestire le API in modo asincorno?
\item
  percher usare httpx.AsyncClient per IncusEngine?
\item
  perche BlueAgent gestisce i container e non mantiene lo stato delle
  challenge?
\item
  motivare le API rests
\item
  perche SQLite
\end{itemize}

La progettazione di BlueAgent risponde alle esigenze operative tipiche
di scenari reali di Blue Teaming, bilanciando fedeltà
dell\textquotesingle ambiente, isolamento dei test ed efficienza delle
risorse.

l\textquotesingle architettura di blueagent ha le seguenti proprieta:

\begin{itemize}
\tightlist
\item
  \textbf{Zero Side-Channel Leakage}: invece di eseguire la verifica
  sull\textquotesingle ambiente di produzione, Blueagent isola la
  verifica su cloni effimeri ed una rete ovn inaccessibile. Dunque un
  partecipante non puo\textquotesingle{} intercettare il payload di test
  (accedendo alle interfacce di rete, guardando i log di systema,
  intercettando le chiamate di sistema)
\item
  \textbf{Nessun impatto sull\textquotesingle ambiente di produzione}:
  Test distruttivi o exploit instabili possono penalizzare
  l\textquotesingle SLA del team, o di compromettere il lavoro del team
  salvato nella macchina.
\item
  \textbf{Fedelta Topologica}: i container mantengo le interfacce come
  definite in origine, gli stessi IP statici e lo stesso gateway, sia in
  produzione che in verifica
\item
  \textbf{Ambiente Stateful Completo:} A differenza di soluzioni come
  kCTF (nsjail/OCI), i container di sistema Incus integrano ambienti con
  \passthrough{\lstinline!systemd!}, log persistenti e architetture
  multi-tier (Web + DB), simulando sistemi reali.
\end{itemize}

inoltre, tecnicamente la clonazione dell\textquotesingle ambiente di
produzione e\textquotesingle{} fattibile poiché:

\begin{itemize}
\tightlist
\item
  clonare un\textquotesingle ambiente con 2 container ubuntu, in un
  filesystem \textbf{ZFS} che offre un meccanismo di \textbf{COW},
  richiede secondo le misurazioni fatte,all\textquotesingle incirca
  \textbf{700 ms}.
\item
  un singolo container di sistema (es: alpine) richiede
  all\textquotesingle incirca 100 RAM. (da misurare). dunque
  e\textquotesingle{} teoricamente possibile instanziare centinaia di
  container all\textquotesingle interno di un singolo host che esegue
  Blueagent.
\item
  la gestione delle risorse morte avviene tramite un demone
  \passthrough{\lstinline!LifecycleManager!} che distrugge le risorse
  orfane dopo 90 secondi di inattivita.
\item
  l\textquotesingle uso di reti OVN permette di replicare la topologia
  di rete in produzione senza generare conflitti nelle tabelle di
  routing nell\textquotesingle host.
\end{itemize}

Voglio precisare qui che ho deciso di utilizzare OVN per gestire la rete
di verifica per evitare l\textquotesingle overhead dovuto ai tunnel
Geneve, a tutto il resto dell\textquotesingle infrastruttura OVN
all\textquotesingle interno dell\textquotesingle ambiente di produzione.

\section{Destinatari e Casi
d\textquotesingle Uso}\label{destinatari-e-casi-duso}

\begin{itemize}
\tightlist
\item
  \textbf{Corsi universitari di sicurezza informatica:} Laboratori
  pratici di patch management e hardening di sistemi Linux completi.
\item
  \textbf{Training aziendale Blue Team / SOC:} Esercitazioni di
  remediation, analisi forense e configurazione di contromisure su
  servizi multi-container.
\item
  \textbf{Competizioni Attack-Defense e Defense-Only:} Gare con decine
  di team che richiedono ambienti persistenti a basso costo
  infrastrutturale.
\end{itemize}

\begin{center}\rule{0.5\linewidth}{0.5pt}\end{center}

\section{Critica e Limiti
Architetturali}\label{critica-e-limiti-architetturali}

Nonostante i vantaggi, l\textquotesingle architettura presenta
compromessi e vincoli strutturali:

\begin{itemize}
\tightlist
\item
  \textbf{UDP Path MTU Blackhole:} La combinazione di tunnel WireGuard
  (\(\text{MTU } 1420\)) e Geneve OVN (\(\text{MTU } 1442\),
  all\textquotesingle interno della \textbf{rete di verifica}) impone
  riduzioni di MTU. Mentre il TCP beneficia del MSS clamping via
  \passthrough{\lstinline!nftables!}, i pacchetti UDP di grandi
  dimensioni con flag \emph{Don\textquotesingle t Fragment}, potrebbero
  essere scartati dal kernel.
\item
  \textbf{Limite scalabilità a 254 team:}
  L\textquotesingle indirizzamento deterministico
  \passthrough{\lstinline!10.k.x.x!} vincola \(k\) a un singolo byte
  (\(0 \le k \le 254\)), quindi ad un massimo di \(\text{255}\) team
\item
  \textbf{Concorrenza e blocco multi-nodo:} Il meccanismo
  \passthrough{\lstinline!TwoTierLock!}
  (\passthrough{\lstinline!asyncio.Lock!} combinato con
  \passthrough{\lstinline!flock!}) è confinato al singolo nodo. Per
  scalare l\textquotesingle architettura su piu nodi bisogna introdurre
  l\textquotesingle uso di un gestore di lock distribuito come Redis.
\end{itemize}
